% !TeX root = _main.tex
% ROLE 1 -- BOI CANH, HE THONG DICH VU, QUY TRINH AS-IS (muc 1--4).
% Nguoi nhan: nhom tu dien. Muc tieu: 4--6 trang noi dung chinh.
% DUOC SUA: file nay, figures/r1_*; nhan \label bat dau bang r1:.
% KHONG SUA: role_2.tex, role_3.tex, _main.tex, _content.tex, analysis/.
% DAU VAO: rubric, PDF Checkpoint 2, CheckPoint1/sections, figures/framework_cp1.png.
% BAN GIAO: pham vi tram do Wq, Resource Map, BPMN AS-IS, nguyen nhan -> bien do.
% HOAN THANH KHI: co BPMN dung ky phap, nguon/nhan AI, khong con TODO role 1.

\section{Bài toán và phạm vi nghiên cứu}
\label{r1:scope}
\subsection{Vấn đề dịch vụ và câu hỏi phân tích}
Dự án xét dịch vụ tiếp nhận triệu chứng trực tuyến, hỗ trợ tổng hợp thông tin,
nhân viên y tế xem xét và hướng dẫn đặt lịch khám khi phù hợp. Báo cáo
Checkpoint 2 đã mô tả hành trình từ khai báo đến xác nhận lịch và chỉ ra điểm
chờ khi hồ sơ được chuyển cho nhân viên y tế~\cite{checkpoint2}.
Báo cáo cuối kỳ tập trung vào một bài toán: thời gian chờ tại trạm tiếp nhận,
phân loại yêu cầu tư vấn và cách phân bổ nhân lực theo khung giờ.

\textbf{Câu hỏi phân tích.} Trong cùng một cấu hình nhu cầu giả định, thay đổi
số điều dưỡng theo giờ và giảm thời gian xử lý hồ sơ có thể giảm thời gian
chờ đến mức nào, với chi phí nhân lực và rủi ro vận hành nào?
Giá trị kỳ vọng là phản hồi kịp thời hơn cho người dùng và sử dụng nguồn lực
phù hợp với thời điểm phát sinh nhu cầu. Kết quả được đánh giá bằng mô phỏng;
chưa phải kết quả triển khai tại cơ sở y tế.

\subsection{Ranh giới hệ thống và ranh giới phép đo}
\begin{table}[htbp]\centering\small
\caption{Phạm vi dịch vụ và phạm vi mô hình định lượng}\label{r1:boundary}
\begin{tabularx}{\linewidth}{@{}p{32mm}Y@{}}
\toprule
Thành phần & Mô tả \\
\midrule
Người sử dụng & Người bệnh hoặc người nhà được ủy quyền khai báo triệu chứng. \\
Tác nhân cung cấp & Điều dưỡng, bác sĩ, nhân viên tiếp nhận và đơn vị vận hành nền tảng. \\
Trigger dịch vụ & Hệ thống nhận mô tả triệu chứng qua ứng dụng. \\
Outcome dịch vụ & Hướng dẫn đã được nhân viên y tế duyệt; xác nhận lịch nếu cần khám và có chỗ. \\
Đầu vào mô hình & Hồ sơ sẵn sàng vào hàng đợi tại trạm điều dưỡng. \\
Đầu ra mô hình & Điều dưỡng hoàn tất xử lý tại trạm; chưa bao gồm toàn bộ thời gian tới khi có lịch khám. \\
Ngoài mô hình & Chẩn đoán, điều trị, thanh toán, công suất bác sĩ và lịch khám phía sau; chất lượng phân loại của AI. \\
\bottomrule
\end{tabularx}
\end{table}
Ứng dụng có thể tiếp nhận thông tin ngoài giờ, nhưng mô phỏng hiện tại chỉ
sinh yêu cầu trong khoảng 07:00--19:00. Việc quy đổi yêu cầu ngoài giờ thành
hàng đợi đầu ngày chưa được mô hình hóa. Nhóm ưu tiên P1 là nhãn của mô hình;
không đồng nhất nó với quy trình xử lý cấp cứu trong toàn bộ dịch vụ.
\todo{1}{Chốt điểm bắt đầu/kết thúc của tác vụ điều dưỡng trên BPMN để khớp với
\texttt{t\_arrival}, \texttt{t\_start}, \texttt{t\_end} ở role 2.}

\section{Mô hình hệ thống dịch vụ}
\label{r1:system}
Báo cáo dùng mô hình hệ thống dịch vụ làm khung chủ đạo để liên kết tác nhân,
nguồn lực, hoạt động, thông tin, quy tắc và kết quả. Góc nhìn đồng tạo giá trị
được dùng để giải thích quan hệ giữa các bên trong cùng khung, theo yêu cầu
một khung phân tích chính của rubric~\cite{rubric}.
\begin{table}[htbp]\centering\small
\caption{Các thành phần của hệ thống và liên hệ với bài toán chờ}
\begin{tabularx}{\linewidth}{@{}p{30mm}Y Y@{}}
\toprule
Thành phần & Biểu hiện trong dịch vụ & Liên hệ với phân tích \\
\midrule
Con người & Người khai báo, điều dưỡng, bác sĩ, tiếp nhận & Số kênh điều dưỡng hoạt động $c_t$. \\
Hoạt động & Thu thập, kiểm tra, phân loại, chuyển tiếp & Thời gian xử lý $S$ và điểm bàn giao. \\
Thông tin & Triệu chứng, tiền sử, mức ưu tiên, trạng thái & Tính đầy đủ có thể làm thay đổi $S$; chưa đo tỷ lệ lỗi. \\
Công nghệ & Biểu mẫu, công cụ tổng hợp, hàng đợi, thông báo & Hỗ trợ nhân viên tiếp nhận thông tin; không suy ra độ chính xác lâm sàng. \\
Quy tắc & Thứ tự xử lý, phân quyền, ngưỡng cảnh báo & FCFS hoặc ưu tiên không chen ngang; mục tiêu thời gian chờ. \\
Kết quả & Hồ sơ được xử lý và chuyển đúng bước & $W_q$, P90, tỷ lệ vượt ngưỡng và chi phí quy ước. \\
\bottomrule
\end{tabularx}
\end{table}
\todo{1}{Bổ sung nguồn học thuật đã đọc cho khung được chọn; dùng thuật ngữ nhất
quán với bài giảng và sơ đồ. Gửi khóa trích dẫn cho người ghép bài.}

\section{Bản đồ nguồn lực và cơ chế đồng tạo giá trị}
\label{r1:resources}
\begin{figure}[htbp]\centering
\includegraphics[width=0.94\linewidth,height=0.48\textheight,keepaspectratio]{figures/framework_cp1.png}
\caption{Bản đồ nguồn lực kế thừa từ Checkpoint 1, cần rà soát theo phạm vi cuối kỳ.}
\label{r1:resource-map}
\source{hình \texttt{framework\_cp1.png} có sẵn trong dự án; lịch sử hỗ trợ AI cần nhóm xác nhận}
\end{figure}
Người sử dụng cung cấp thông tin và phối hợp trả lời; nhân viên y tế đóng góp
tri thức và thời gian; cơ sở khám cung cấp năng lực tiếp nhận; đơn vị vận hành
cung cấp hạ tầng và cơ chế phối hợp. Giá trị xuất hiện khi các nguồn lực được
kết nối thành một yêu cầu có đủ thông tin để xử lý và chuyển tiếp.

Chậm phản hồi có thể phát sinh khi thời gian điều dưỡng không khớp nhu cầu
theo giờ. Thông tin thiếu làm tăng công việc bổ sung và có thể kéo dài xử lý.
Hai cơ chế này tương ứng với hai biến quyết định trong mô hình: $c_t$ và
$E[S]$. Tuy nhiên, mức giảm $E[S]$ nhờ công cụ hỗ trợ hiện là giả định;
bản đồ nguồn lực tự nó không chứng minh hiệu quả của công cụ.
\todo{1}{Hoàn thiện Resource Map: đánh dấu ranh giới, 3--5 nhóm nguồn lực,
mũi tên trao đổi, chú giải và lợi ích/chi phí/rủi ro. Nếu vẽ lại, lưu
\texttt{figures/r1\_resource\_map.pdf} và thay hình trên.}

\section{Quy trình hiện tại và chẩn đoán điểm nghẽn}
\label{r1:asis}
\subsection{Journey và Service Blueprint AS-IS}
Hành trình tham chiếu gồm khai báo triệu chứng, hỏi bổ sung, chờ duyệt,
nhận hướng dẫn, chọn lịch và xác nhận lịch~\cite{checkpoint2}. Hai điểm cần
theo dõi là thiếu thông tin khi khai báo và không biết tiến độ trong lúc chờ.
Blueprint cần nối các hành động này với xử lý phía sau, thay vì chỉ liệt kê
các màn hình ứng dụng.
\begin{table}[htbp]\centering\small
\caption{Tóm tắt các lớp của blueprint tại điểm chờ}
\begin{tabularx}{\linewidth}{@{}p{36mm}Y@{}}
\toprule
Lớp & Nội dung ở bước chuyển hồ sơ cho điều dưỡng \\
\midrule
Evidence & Mã hồ sơ, trạng thái tiếp nhận và kết quả phản hồi. \\
Customer actions & Gửi hồ sơ, theo dõi, cung cấp thông tin bổ sung. \\
Frontstage & Hiển thị trạng thái; nhân viên tương tác khi cần bổ sung. \\
Backstage & Kiểm tra hồ sơ, xếp hàng, điều dưỡng tiếp nhận và chuyển tiếp. \\
Support processes & Lịch trực, phân quyền, lưu dấu thời gian và thông báo. \\
\bottomrule
\end{tabularx}
\end{table}
\todo{1}{Bổ sung hoặc chuyển Journey/Blueprint từ PDF mẫu; chỉ rõ line of
interaction, line of visibility và ít nhất hai pain points. Không gọi bảng
tóm tắt này là sơ đồ blueprint hoàn chỉnh.}

\subsection{BPMN AS-IS và ngoại lệ}
\IfFileExists{figures/r1_bpmn_asis.pdf}{%
\begin{figure}[htbp]\centering
\includegraphics[width=\linewidth,height=0.65\textheight,keepaspectratio]{figures/r1_bpmn_asis.pdf}
\caption{BPMN AS-IS của quy trình tiếp nhận và phân loại.}\label{r1:bpmn}
\end{figure}
}{\placeholder{BPMN AS-IS -- role 1}{Xuất sơ đồ thành
\texttt{figures/r1\_bpmn\_asis.pdf}; báo cáo sẽ tự chèn khi file tồn tại.
Cần có start/end, lane/pool, task, gateway khi có rẽ nhánh, nhánh thiếu
thông tin và điểm bàn giao vào hàng đợi.}}
\todo{1}{Đối chiếu \texttt{assests/diagram.bpmn} và các sơ đồ checkpoint;
chọn đúng AS-IS. Kiểm tra sequence flow trong pool, message flow giữa pool,
nhánh ngoại lệ và ghi nguồn/AI ngay dưới sơ đồ.}

\subsection{Từ điểm nghẽn đến biến đo}
\begin{table}[htbp]\centering\small
\caption{Các giả thuyết chẩn đoán cần đối chiếu với dữ liệu}
\begin{tabularx}{\linewidth}{@{}Y Y Y@{}}
\toprule
Điểm trên quy trình & Cơ chế giả định & Biến kiểm tra \\
\midrule
Chờ điều dưỡng nhận hồ sơ & Công suất cố định trong khi nhu cầu biến thiên & $\lambda_t$, $c_t$, $\rho_t$, $W_q$. \\
Yêu cầu bổ sung thông tin & Thiếu dữ liệu làm kéo dài thao tác & $E[S]$, CV; mô hình chưa tách riêng vòng rework. \\
Hồ sơ ưu tiên chờ cùng hàng & FCFS không phản ánh mức ưu tiên & $W_q$ theo P1/P2/P3. \\
Chuyển tiếp sang lịch khám & Có thể có điểm nghẽn sau trạm điều dưỡng & Cần dữ liệu riêng; ngoài mô phỏng hiện tại. \\
\bottomrule
\end{tabularx}
\end{table}
